\section{Interpretation of the experiments}
\label{sec:experimental-interpretation}

The main experiment finds a clear structure.  One quadratic controls the
interior of the cross-wedge.  Three affine boundary lines contain extra
quotient directions.  The same failures remain when the Fourier cutoff is
increased.  They are therefore separate filtration cases, not truncation
errors.

The proofs explain several parts of this pattern for all admissible
parameters.  The first-Hasse-jet formula explains why the $D=2$ block fails
in the new wedge.  Exact filtration shows that the three boundary directions
survive in the quotient.  The perfect Hasse-quotient pairing and theta
skew-adjointness connect the two obstructions.  These theorems do not give
every boundary coefficient.  The missing steps are the factorial block
reduction in Conjecture~\ref{conj:factorial-block-reduction} and two
triangular containments for the common endpoint core.  The pre-low-point and
reset jet terms are now evaluated uniformly and lie below the relevant
cutoffs.

The evidence is strongest on the third boundary line.  The one-block shape
was first observed in $78$ cells.  A separate blind audit tested $652$ cells.
A second audit used $186$ ordinary forms in a two-dimensional cusp-form
space.  The two later audits made $2{,}512$ tests.  Every test returned the
same form-independent factorial
scalar.  The formula is a precise and well-tested conjecture.  The finite
tests do not replace the missing symbolic proof.

There are three limits.  First, the main cross-wedge table ends
at $p=43$, although the separate third-line test reaches $p=79$ and the
symbolic index audit reaches $p=101$.  Second, the broad record uses
level-one ordinary forms; supersingular and higher-level behavior is not
tested.  Third, the two triangular coefficient identities have exact finite
support, proved endpoint independence, and proved low-point jet reductions,
but the uniform containment of the remaining core series is still open.
These limits give clear next tests and proof problems.
